\conclusion

В ходе работы разработана гибридная каскадная система автоматизированного обогащения товарных данных для интернет-магазина. Система реализована в виде программного комплекса с демонстрационным стендом, проверена на выборке без пересечения брендов и удовлетворяет шести функциональным, пяти нефункциональным и четырём архитектурным требованиям ТЗ.

Достигнутые показатели превосходят целевые ориентиры ТЗ (≥ 85 \% точности, ≥ 2,5× сокращения стоимости): сквозная точность 92,8 \%, стоимость сокращена в 720 раз относительно прямого вызова коммерческой LLM. Каскад без запасного слоя покрывает 96,7 \% ячеек со средней точностью 90,5 \%; в LLM направляется лишь 3,3 \% сложных ячеек.

\subsection*{Сводные выводы по главам}

\textbf{Глава 1} систематизирует литературу и обосновывает противоречие. PIM-системы (Akeneo \cite{ref_5}, Pimcore \cite{ref_6}, Salsify \cite{ref_7}) ориентированы на хранение и распределение, а появляющиеся модули автоматического обогащения — обёртки над одной LLM, не отвечающие требованиям больших каталогов. Академические подходы разделены: специализированные системы извлечения признаков (TXtract \cite{ref_23}, MAVE \cite{ref_24}, OpenTag \cite{ref_25}) обеспечивают качество без учёта экономики; каскады LLM (FrugalGPT \cite{ref_9}, AutoMix \cite{ref_10}, Hybrid LLM \cite{ref_11}, RouterBench \cite{ref_12}) дают инструменты управления стоимостью, но не интегрируют классическое ML. Из противоречия выведена гипотеза о каскаде разнородных слоёв и ТЗ.

\textbf{Глава 2.} Задача формализована как поатрибутное предсказание со схемой, допускающей явный отказ. Логика — каскад из пяти слоёв с последовательной фильтрацией: предсказание Слой 2 принимается только при одновременном выполнении $P_{\text{ML}} \geq \tau_k$ и $P_B \geq \theta_k$ — формализация ансамблевой оценки уверенности как AND-условия. Архитектурный выбор каждой компоненты (SBERT, XGBoost, pgmpy, Go-шлюз + Python-ML, открытая gpt-oss-120b) обоснован эксплуатационно, методологически и стратегически.

\textbf{Глава 3.} Каскад реализован как трёхкомпонентный комплекс (Go-шлюз, Python/FastAPI ML-сервис, статический веб-интерфейс); состав — 24 поатрибутных XGBoost (8/7/9 для pasta/chocolate/cheeses), три байесовские сети и Слой 0. На 4350 ячейках без пересечения брендов сквозная точность каскад + Gemini 2.5 Flash — 92,8 \%, на 9,0 п. п. выше прямого Sonnet 4.5 при стоимости в 720 раз ниже; в сценарии локального развёртывания каскад + gpt-oss-120b — 92,3 \% (на 23,0 п. п. выше самостоятельного использования). Слой 1 закрывает 18–23 \% pasta/chocolate с точностью 87–90 \% и 3 \% cheeses со 100 \% precision; Слой 2 — 72–94 \% покрытия при 93,5–98,6 \% точности; байесовский валидатор работает селективно на chocolate/contains\_nuts; LLM добавляет 2,3 п. п. при 3,3 \% обращений. Зафиксированная до получения результатов H1 о превосходстве обучаемого маршрутизатора отклонена. Цикл «аудит — правки — переобучение — измерение» воспроизведён 3/3 с приростом +18,2, +16,3 и +14,4 п. п.

\textbf{Глава 4} переводит результаты в плоскость применения: три сценария по выбору слоя 4 (gemini-flash, Sonnet/GPT-4o, локально развёрнутая gpt-oss-120b); три ролевые модели (контент-менеджер, системный аналитик, инженер ML). Сводный эффект — сокращение стоимости в 720 раз и ручной нагрузки в 10–20 раз при точности 92,8 \% с интервалом Вильсона [92,0; 93,6].

\subsection*{Закрытие формулировок утверждённой темы}

Утверждённая тема работы содержит восемь содержательных формулировок. Каждой из них в работе соответствует конкретный раздел, а также архитектурный или экспериментальный компонент разработки. Соответствие представлено в таблице 5.1.

\begin{xltabular}{\textwidth}{|X|X|X|}
\caption{Соответствие формулировок утверждённой темы и разделов работы} \\ \hline
    Формулировка из обоснования актуальности темы & Реализация в работе & Раздел \\ \hline
\endfirsthead
\multicolumn{3}{l}{\textit{Продолжение таблицы}} \\ \hline
    Формулировка из обоснования актуальности темы & Реализация в работе & Раздел \\ \hline
\endhead
    Гибридная каскадная система & Пятислойная архитектура каскада (Слой 0 + четыре содержательных слоя) & 2.2, 3.1 \\ \hline
    Семантические возможности генеративных моделей & Запасной слой 4 на основе большой языковой модели & 3.1.1 \\ \hline
    Предсказательная сила классического машинного обучения & Слой 2 — поатрибутные классификаторы XGBoost на многоязычных векторных представлениях SBERT & 3.1.1 \\ \hline
    Ансамблевая оценка уверенности моделей & Многоуровневая фильтрация по уверенности как последовательное конъюнктивное условие на калиброванные пороги слоёв & 2.1 \\ \hline
    Отбраковка фактических ошибок до их попадания в каталог & Слой 3 — байесовский валидатор плотностей пар «бренд × значение» & 3.1.1, 3.3.4 \\ \hline
    Скрытые статистические закономерности каталога & Направленный ациклический граф атрибутов, обучаемый алгоритмом Hill Climb с критерием BIC & 3.1.1 \\ \hline
    Разнородные источники данных (мультимодальные и многоканальные) & OCR-извлечение из изображений Open Food Facts и эмбеддинги CLIP на этапе обучения; нечёткое сопоставление с базой USDA Branded Foods для производных атрибутов & 3.2.5 \\ \hline
    Интернет-магазин (партнёр — каталог — витрина) & Демонстрационный программный комплекс с REST-интерфейсом обработки партнёрского ввода и руководством пользователя & 1.1, 3.1.2, 4.1 \\ \hline
\end{xltabular}

Таким образом, все восемь содержательных формулировок утверждённой темы имеют прямое отображение в разработанной системе и подтверждены экспериментально либо архитектурно. Наиболее нагруженные формулировки — отбраковка фактических ошибок и использование скрытых статистических закономерностей каталога — закрыты слоем байесовского валидатора. Для него в работе пересмотрена роль: вместо самостоятельного классификатора сеть используется как валидатор плотностей.

\subsection*{Научная новизна}

Научная новизна работы складывается из трёх содержательных пунктов.

\begin{enumerate}[arabic]
    \item Доминирование каскадной архитектуры по Парето над всеми безусловными конфигурациями большой языковой модели на задаче обогащения товарных атрибутов. Одновременно достигается превышение точности на 9,0 процентных пункта и снижение стоимости в 720 раз по сравнению с прямым обращением к Claude Sonnet 4.5. Результат подтверждён на проверочной выборке без пересечения брендов из 4350 ячеек и воспроизведён на пяти различных моделях запасного слоя. Это означает, что эффект сохраняется при смене поставщика генеративной модели.
    \item Сценарий развёртывания на собственных мощностях с открытой большой языковой моделью. Каскад с открытой моделью gpt-oss-120b достигает 92,3 \% сквозной точности — на 23,0 процентных пункта выше, чем самостоятельное применение той же модели. Результат демонстрирует способность каскадной композиции компенсировать слабость малых открытых моделей через специализированные слои. Кроме того, он формирует методологическую основу для локального развёртывания системы без передачи партнёрских данных во внешние коммерческие сервисы.
    \item Зафиксированный до получения результатов отрицательный результат для обучаемого стоимостно-чувствительного маршрутизатора и кросс-доменная репликация цикла «аудит — исправления — переобучение — измерение». Гипотеза о превосходстве обучаемого XGBoost-маршрутизатора над статическим правилом отклонена на трёх бюджетах с поправкой Бонферрони. Этот результат согласуется с выводами бенчмарка RouterBench \cite{ref_12}. Цикл аудита эталонной разметки воспроизведён на трёх категориях из трёх с количественно измеримым приростом сквозной точности.
\end{enumerate}

\subsection*{Ограничения работы}

Результаты ограничены рядом обстоятельств, добросовестно фиксируемых для формирования перспектив. Во-первых, область определения ограничена тремя категориями пищевой продукции; переносимость на другие домены не валидирована. Во-вторых, эталон построен сочетанием слабого надзора и согласия трёх LLM со слепой проверкой Opus 4 — аудит-эталон второго уровня, ручная экспертная разметка не выполнялась. В-третьих, мультимодальные сигналы OCR и CLIP используются только на обучении; в выводе работают только текстовые поля. В-четвёртых, проверка ограничена пятью европейскими языками, устойчивость к иным письменностям не валидирована. В-пятых, предварительная регистрация ограничена H1; показатели 3.3.2 — post-hoc анализ.

Селективный характер байесовского валидатора (значимый прирост точности отбраковки на 3 парах из 22) ограничением не считается, а признаётся архитектурным свойством. В финальной конфигурации валидатор включён точечно на chocolate/contains\_nuts, где даёт +6,6 п. п. точности при росте расходов на 5,3 п. п. (3.3.4). Принцип «изменение роли компонента вместо его удаления» сформулирован в 3.1.1 и 3.3.4 как самостоятельный методологический вывод.

\subsection*{Перспективы развития темы}

Полученные результаты и зафиксированные ограничения определяют ряд перспективных направлений развития темы.

\begin{enumerate}[arabic]
    \item Расширение области определения до шести категорий — напитки, зерновые завтраки и косметика. Для этих категорий подготовлены схемы атрибутов и серебряная разметка в файлах \texttt{beverages\_stratified\_silver\_standard.parquet}, \texttt{cereals\_stratified\_silver\_standard.parquet} и \texttt{cosmetics\_stratified\_silver\_standard.parquet}. Основной блокер расширения — построение аудит-эталона уровня итоговой эталонной разметки v2 для каждой категории и обучение поатрибутных классификаторов Слой 2. Методология аудита и переобучения предположительно переносится на новые категории без принципиальных изменений.
    \item Регулярная переаттестация изотонической калибровки на следующих итерациях эталонной разметки. Эффект схемы перекрёстной калибровки экспериментально проверен и реализован в коде системы (3.3.3.3, флаг \texttt{--calibration-method isotonic} в \texttt{src/pipeline/ml/train.py}). Средняя ECE снижается с 0,070 до 0,043, что даёт улучшение на 2,78 п. п. в среднем и прирост на 18 из 22 атрибутов. Перспектива заключается в переаттестации после расширения эталонной разметки до большего числа атрибутов и категорий. Также целесообразно добавить регрессионный тест на ухудшение ECE более чем на 0,01 как условие принятия следующего цикла переобучения.
    \item Введение мультимодального вывода в производственный режим — расширение прикладного интерфейса на приём изображений и добавление OCR/CLIP в реальном времени. Такой шаг требует пересмотра контракта взаимодействия с партнёром.
    \item Запуск пилотного активного обучения на ячейках расхождения каскада и запасного слоя \cite{ref_46, ref_39}. Логика построения очереди ячеек, где предсказание слоя 2 расходится с ответом запасного слоя, реализована в \texttt{src/experiments/active\_learning\_pilot.py}. Скрипт выполняет выборку 5 000 кодов на категорию, отбор 150 неуверенных и 150 случайных контрольных, а также повторное обучение с замером прироста на отложенной выборке. Перспектива заключается в проведении пилота в полном объёме на трёх категориях и в измерении прироста точности слоя 2 относительно случайной разметки той же стоимости.
    \item Использование базы USDA Branded Foods и аналогов как структурированного второго источника для производных атрибутов с систематизацией политики приоритетов между партнёрским вводом и вторичными каталогами.
    \item Локализация на русский язык и российские каталоги — разработка русскоязычных регулярных выражений слоя 1 и дообучение слоя 2 на каталогах российских ритейлеров. Поскольку SBERT поддерживает русский в едином векторном пространстве, трудоёмкость ограничена слоями 1–2.
    \item Перенос архитектуры в непродовольственные домены, начиная с электроники. Для смартфонов из каталога PhoneDB подготовлены схема атрибутов, скрипт построения серебряного стандарта и демонстрация холодного старта. Полноценное измерение точности требует построения аудит-эталона уровня итоговой эталонной разметки v2 и оценки без пересечения брендов. Электроника сохранена как контрольный домен проверки переносимости и в основную таблицу 3.3 не входит.
\end{enumerate}

Таким образом, разработанная гибридная каскадная система демонстрирует возможность одновременно достигать высокой точности и низкой стоимости обогащения товарных каталогов. Этот результат обеспечивается за счёт стратифицированного распределения нагрузки между разнородными по природе слоями обработки. Полученные результаты подтверждают выдвинутую в первой главе работы архитектурную гипотезу о каскадной композиции и закрывают все восемь содержательных формулировок утверждённой темы.
